% =====================
\section{リスク管理}

担当B/Cに関わるリスクを表\ref{tab:risks}に示す．

\begin{table}[H]
  \centering
  \begin{tabularx}{\linewidth}{lXXl}
    \toprule
    リスク & 影響 & 対応策 & 担当 \\
    \midrule
    赤外線センサー誤検出 &
      テンポ推定が不安定になり演奏がずれる &
      複数サンプルの移動平均でノイズ除去．誤検出時は直前テンポを保持する & B \\
    フォトTrの光量不足 &
      拍を検出できず演奏が停止する &
      LED輝度・検出距離を単体テストで調整．閾値をTPM-2で管理する & A/B \\
    シリアル通信の遅延増大 &
      楽器間の発音タイミングがずれる &
      遅延補正値（PC1c）で補正．結合テストC1で実測・確定する & B/C \\
    楽譜ポインタのずれ &
      フェルマータ後に音符が正しく進まない &
      フラグ管理方式（\texttt{fermata}フラグ）で対応．ブロッキング処理は使用しない & C \\
    テンポ変化時の収束遅延 &
      2周目（104\,bpm）への移行が遅れ演奏が乱れる &
      テンポ変化検出時は$\alpha=0.5$へ切り替えて素早く収束させる & B \\
    \bottomrule
  \end{tabularx}
  \caption{リスク一覧（担当B/C）}
  \label{tab:risks}
\end{table}
